% Irodov Part III §3.3 problems 101--112, user PDF pp.114--116,
% printed pp.105--107. Capacitor figure 3.16--3.17 inspected.
\electrocapproblem{101}{105}
半径为 $R_1$ 的孤立金属球外表紧贴同心介质壳，
介质相对介电常数为 $\varepsilon$，外半径 $R_2$；
其外部真空。求金属球相对于无穷远的电容。
\exsolution
设球总自由电荷为 $q$。高斯定律给
$D_r=q/(4\pi r^2)$，故
$R_1<r<R_2$ 时
$E_r=q/(4\pi\varepsilon_0\varepsilon r^2)$，
$r>R_2$ 时
$E_r=q/(4\pi\varepsilon_0r^2)$。
球电势
\[
\varphi(R_1)=\frac{q}{4\pi\varepsilon_0}
\left[\frac1\varepsilon
\left(\frac1{R_1}-\frac1{R_2}\right)
+\frac1{R_2}\right].
\]
由 $C=q/\varphi(R_1)$：
\[
\boxed{C=
\frac{4\pi\varepsilon_0\varepsilon R_1}
{1+(\varepsilon-1)R_1/R_2}}.
\]
\excheck $R_2=R_1$ 时退化为真空导体球电容
$4\pi\varepsilon_0R_1$；
$R_2\to\infty$ 时得到
$4\pi\varepsilon_0\varepsilon R_1$。

\electrocapproblem{102}{105}
两个空气平板电容器电容均为 $C$，
串联接在电动势 $\mathcal E$ 不变的电池两端。
随后将其中一个电容器完全填充
相对介电常数为 $\varepsilon$ 的介质。
问该电容器的电场减小多少倍，
并求经电池流过的电荷量。
\exsolution
原串联等效电容 $C_{\rm eq,0}=C/2$，
串联电荷 $Q_0=C\mathcal E/2$，
每个电容器各承压 $\mathcal E/2$。
填充后该电容器电容为 $\varepsilon C$，
新串联电容为
$C_{\rm eq}=\varepsilon C/(1+\varepsilon)$。
故它的新电压
$V_{\rm diel}=Q/(\varepsilon C)
=\mathcal E/(1+\varepsilon)$。
同一板距下 $E\propto V$，因此
\[
\boxed{\frac{E_{\rm new}}{E_{\rm old}}
=\frac2{\varepsilon+1}},
\quad
\text{原电场减小 }(\varepsilon+1)/2\text{ 倍}.
\]
通过电池的电荷量等于两次串联电荷差：
\[
\boxed{\Delta Q=
C\mathcal E\left[\frac{\varepsilon}{1+\varepsilon}
-\frac12\right]
=\frac{C\mathcal E(\varepsilon-1)}
{2(\varepsilon+1)}}.
\]
\excheck 电源连接未断开，必须保持
$\mathcal E$ 不变，而不能错误地保持串联电荷不变。

\electrocapproblem{103}{105--106}
平板电容器板面积为 $S$，间距由相邻
两层介质填满；厚度 $d_1,d_2$，
相对介电常数 $\varepsilon_1,\varepsilon_2$。
求电容及两层界面处束缚面电荷密度
（场从1层指向2层，外加电压 $V$）。
\exsolution
无自由界面电荷，法向位移量恒定 $D$，
两层 $E_i=D/(\varepsilon_0\varepsilon_i)$。
电压
$V=D[d_1/\varepsilon_1+d_2/\varepsilon_2]/\varepsilon_0$。
因此
\[
\boxed{C=
\frac{\varepsilon_0S}{d_1/\varepsilon_1+d_2/\varepsilon_2}}.
\]
两层极化分别为
$P_i=D(1-1/\varepsilon_i)$。
在界面上定义沿1层指向2层为正法向，
来自介质1的束缚面电荷为 $+P_1$，
来自介质2的为 $-P_2$，
相加得
\[
\sigma_b=P_1-P_2
=D\left(\frac1{\varepsilon_2}-\frac1{\varepsilon_1}\right).
\]
代入
$D=\varepsilon_0V/(d_1/\varepsilon_1+d_2/\varepsilon_2)$：
\[
\boxed{\sigma_b=
\frac{\varepsilon_0V(\varepsilon_1-\varepsilon_2)}
{\varepsilon_1d_2+\varepsilon_2d_1}}.
\]
\excheck 这是\emph{同一界面两侧束缚电荷之和}，
不是单取其中一层的表面极化电荷。

\electrocapproblem{104}{106}
平板电容器面积 $S$，板距 $d$，
介质相对介电常数沿电场方向线性地
由 $\varepsilon_1$ 增至 $\varepsilon_2>\varepsilon_1$。
求电容及带电 $q$ 时束缚体电荷密度
关于局部 $\varepsilon$ 的关系。
\exsolution
没有自由体电荷，$D=q/S$ 恒定，
$E(x)=D/[\varepsilon_0\varepsilon(x)]$。
令
$\varepsilon(x)=\varepsilon_1+
(\varepsilon_2-\varepsilon_1)x/d$，
则
\[
V=\frac{q}{\varepsilon_0S}
\int_0^d\frac{\mathrm dx}{\varepsilon(x)}
=\frac{qd}{\varepsilon_0S(\varepsilon_2-\varepsilon_1)}
\ln\frac{\varepsilon_2}{\varepsilon_1}.
\]
因此
\[
\boxed{C=
\frac{\varepsilon_0S(\varepsilon_2-\varepsilon_1)}
{d\ln(\varepsilon_2/\varepsilon_1)}}.
\]
$P=D(1-1/\varepsilon)$，
$\rho_b=-\mathrm dP/\mathrm dx$，故
\[
\boxed{\rho_b(x)=
-\frac{q}{S}\frac{\varepsilon_2-\varepsilon_1}
{d\,\varepsilon(x)^2}}.
\]
\excheck 若 $\varepsilon_2\to\varepsilon_1=\varepsilon$，
洛必达法则给 $C\to\varepsilon\varepsilon_0S/d$、
$\rho_b\to0$。

\electrocapproblem{105}{106}
半径为 $R_1,R_2$ 的同心球形电容器，
$R_2>R_1$。填充介质相对介电常数随半径
变化为 $\varepsilon(r)=a/r$，
其中 $a$ 为具有长度量纲的常数。
求电容。
\exsolution
径向位移
$D_r=q/(4\pi r^2)$，
电场
$E_r=D_r/[\varepsilon_0(a/r)]
=q/(4\pi\varepsilon_0 a r)$。
电压
\[
V=\int_{R_1}^{R_2}E_r\,\mathrm dr
=\frac{q}{4\pi\varepsilon_0a}
\ln\frac{R_2}{R_1}.
\]
因此
\[
\boxed{C=\frac{4\pi\varepsilon_0a}
{\ln(R_2/R_1)}}.
\]
\excheck 因 $\varepsilon=a/r$ 为无量纲，
$a$ 必须具有长度量纲。

\electrocapproblem{106}{106}
同轴圆柱电容器两层介质分别占
$R_1<r<R_2$ 与 $R_2<r<R_3$，
介电常数为 $\varepsilon_1,\varepsilon_2$，
击穿场强为 $E_{1m},E_{2m}$。
求增大电压时两层同时达到各自击穿场的条件。
\exsolution
每单位轴向长度自由电荷量记为 $\lambda$。
由柱面高斯定律
$D_r=\lambda/(2\pi r)$，
因此第 $i$ 层
$E_i(r)=|\lambda|/(2\pi\varepsilon_0\varepsilon_i r)$。
场强随半径增大而减小，
所以第1层最大场在 $R_1$，
第2层最大场在 $R_2$。
若两层同步击穿：
\[
\frac{|\lambda|}{2\pi\varepsilon_0\varepsilon_1R_1}
=E_{1m},\quad
\frac{|\lambda|}{2\pi\varepsilon_0\varepsilon_2R_2}
=E_{2m}.
\]
消去 $\lambda$：
\[
\boxed{\varepsilon_1R_1E_{1m}
=\varepsilon_2R_2E_{2m}}.
\]
\excheck 不能用相同的径向位置比较不同层击穿，
每层的危险位置是其各自内表面。

\electrocapproblem{107}{106}
图3.16的双层同轴圆柱电容器，
半径依次为 $R_1<R_2<R_3$，
介质参数同上。
若
$\varepsilon_1R_1E_{1m}<
\varepsilon_2R_2E_{2m}$，
求整体击穿电压。
\exsolution
由于第1层允许的线电荷阈值更小，
先在 $r=R_1$ 达到 $E_{1m}$，
此时
$|\lambda|=2\pi\varepsilon_0\varepsilon_1R_1E_{1m}$。
电压由两层径向积分：
\[
V=\frac{|\lambda|}{2\pi\varepsilon_0}
\left[\frac1{\varepsilon_1}\ln\frac{R_2}{R_1}
+\frac1{\varepsilon_2}\ln\frac{R_3}{R_2}\right].
\]
代入阈值：
\[
\boxed{V_{\rm breakdown}=
R_1E_{1m}
\left[\ln\frac{R_2}{R_1}
+\frac{\varepsilon_1}{\varepsilon_2}\ln\frac{R_3}{R_2}\right]}.
\]
\excheck 条件为严格小于，因此先击穿第1层；
同号等号时两层同时击穿。

\electrocapproblem{108}{106}
空气中两根等长平行细金属导线，
截面半径均为 $a$，轴间距 $b\gg a$。
求两导线\emph{每单位长度的互电容}。
\exsolution
设两线单位长度带自由电荷
$+\lambda,-\lambda$。
二维线电荷势
$\varphi\propto\ln r$。
以 $a\ll b$ 的近似在正线表面取
距正线 $a$、负线 $b$，
同理负线反向，得到电压
\[
V\simeq\frac{\lambda}{\pi\varepsilon_0}\ln\frac ba.
\]
故单位长度电容
\[
\boxed{C'=\frac{\lambda}{V}
\simeq\frac{\pi\varepsilon_0}{\ln(b/a)}}.
\]
\excheck 更精确的两圆柱公式
$C'=\pi\varepsilon_0/
\operatorname{arcosh}(b/(2a))$，
在 $b\gg a$ 化成上述结果。

\electrocapproblem{109}{106}
一根半径 $a$ 的长金属导线，
平行于无限金属导体平面，
轴线距平面 $b\gg a$。
求两者的每单位长度互电容。
\exsolution
用导体平面镜像法，
真实导线带线电荷 $+\lambda$，
其镜像为距真线 $2b$ 的 $-\lambda$。
真线与镜像的电势差约为
$(\lambda/\pi\varepsilon_0)\ln(2b/a)$，
但平面为对称零电势位置，
真线到平面的电压为其中一半：
\[
V\simeq\frac{\lambda}{2\pi\varepsilon_0}\ln\frac{2b}{a}.
\]
故
\[
\boxed{C'=\frac{\lambda}{V}
\simeq\frac{2\pi\varepsilon_0}{\ln(2b/a)}}.
\]
\excheck 对无限导体平面并未采用孤立双线的
完整电压，否则会误少一倍。

\electrocapproblem{110}{106}
两个半径同为 $a$ 的导体球，
球心距 $b\gg a$，
均浸于相对介电常数为 $\varepsilon$ 的均匀介质中。
求双球系统的互电容。
\exsolution
设两球带电 $+q,-q$，
在 $b\gg a$ 下近似各自球面电荷均匀。
正球电势约
$\varphi_+=kq/(\varepsilon a)-kq/(\varepsilon b)$，
负球电势相反。
因此
\[
V=\varphi_+-\varphi_-=
\frac{q}{2\pi\varepsilon_0\varepsilon}
\left(\frac1a-\frac1b\right).
\]
电容
\[
C=\frac qV=
\frac{2\pi\varepsilon_0\varepsilon}
{1/a-1/b}
\simeq\boxed{2\pi\varepsilon_0\varepsilon a}
\quad(b\gg a).
\]
\excheck $b$ 有限时球面电荷并非严格均匀；
上述 $1/b$ 项仅用于说明首阶近似。

\electrocapproblem{111}{106}
半径为 $a$ 的金属球，
球心距无限导体平面为 $\ell\gg a$。
求球—平面的互电容。
\exsolution
镜像法使金属球附近的电势
近似为球自身带电 $q$ 的电势
减去镜像 $-q$ 在距 $2\ell$ 处的电势：
\[
V\simeq\frac{q}{4\pi\varepsilon_0}
\left(\frac1a-\frac1{2\ell}\right).
\]
因此
\[
C\simeq\frac{4\pi\varepsilon_0}
{1/a-1/(2\ell)}
\longrightarrow
\boxed{4\pi\varepsilon_0a}\quad(\ell/a\to\infty).
\]
\excheck 球面电荷的畸变为更高阶修正，
本题只要求 $\ell\gg a$ 的主阶。

\electrocapproblem{112}{106--107}
原书图3.17(a)、(b)为两种电容器连接网络，
求 A、B 端等效电容，按图中标号
(a) $C_1,C_2,C_3$，(b) 每只电容均为 $C$。
\exsolution
(a) 图3.17(a) 中三只电容器各连接同一对
等势节点 A、B；它们两端的电势差相等，
故各自电荷 $Q_i=C_iV$，
总电荷 $Q=\sum_iQ_i$，
因此
\[
\boxed{C_{AB}^{(a)}=C_1+C_2+C_3}.
\]
(b) 图3.17(b) 的四个边电容构成
对称桥式结构；沿对称中心线的节点电势相等，
可将桥的中间电容视为无压差。
剩余两条支路各有两个相等电容串联，
每支电容 $C/2$，两支并联：
\[
\boxed{C_{AB}^{(b)}
=\frac C2+\frac C2=C}.
\]
\excheck 基于原书图3.17网络连线，
不是将四只电容全串联或全并联；
按 $C_i$ 同等电容极限可复核。
